% Generated draft: requires mathematical and source review.
\documentclass[11pt]{article}
\usepackage[margin=25mm]{geometry}
\usepackage{fontspec}
\setmainfont{Latin Modern Roman}
\usepackage{amsmath,amssymb,mathtools,unicode-math}
\setmathfont{Latin Modern Math}
\usepackage{enumitem,needspace,xurl,hyperref}
\setlength{\parindent}{0pt}
\setlength{\parskip}{5pt}
\setlength{\emergencystretch}{3em}
\newcommand{\N}{\mathbb{N}}
\newcommand{\Z}{\mathbb{Z}}
\newcommand{\Q}{\mathbb{Q}}
\newcommand{\R}{\mathbb{R}}
\newcommand{\C}{\mathbb{C}}
\newcommand{\F}{\mathbb{F}}
\newcommand{\A}{\mathbb{A}}
\newcommand{\PP}{\mathbb{P}}
\newcommand{\E}{\mathbb{E}}
\newcommand{\eps}{\varepsilon}
\newcommand{\dd}{\,\mathrm{d}}
\begin{document}
\section*{UnsolvedMath ID: 1100802}
\textbf{Source problem code:} catalogue entry
\subsection{Definitions and mathematical statement}
An $\mathbb R$-tree is a metric space in which any two points are joined by a unique topological arc and that arc is isometric to a real interval. A simplicial tree is a connected graph with no cycles; its geometric realization has the path metric with unit edge lengths. Here actions on $\mathbb R$-trees are by isometries, and actions on simplicial trees are by graph automorphisms. An action on a tree $T$ is nontrivial when
\[
\{x\in T:gx=x\text{ for every }g\in G\}=\varnothing.
\]
For simplicial trees, fixed points are understood in the geometric realization; one may equivalently subdivide to remove inversions.

\textbf{Question (Shalen).} For every finitely presented group $G$, if there exist an $\mathbb R$-tree $T$ and an isometric action $G\curvearrowright T$ with no global fixed point, must there exist a simplicial tree $S$ and an action $G\curvearrowright S$ with no global fixed point? The simplicial action need not arise by discretizing the given action, and no hypothesis on faithfulness or arc stabilizers is imposed.

The supplied summary reports an affirmative claim in Dunwoody's 2022 arXiv version, with explicit reservations: no journal publication or independent confirmation was found there, the proof was not checked line by line, and earlier versions asserted the opposite answer using the Higman group. The record therefore classifies the problem as partial progress rather than an independently established resolution. For the finitely generated setting it reports counterexamples with property (FA), meaning every action on a simplicial tree without inversions has a global fixed point, but with nontrivial $\mathbb R$-tree actions even having finite arc stabilizers.
\subsection{Short English statement}
Does the existence of any fixed-point-free $\mathbb R$-tree action of a finitely presented group force some fixed-point-free simplicial-tree action? The supplied affirmative preprint claim carries substantial verification caveats; the finite-presentation assumption is essential to the reported scope.
\subsection{Sources}
\begin{itemize}
\item M. Bestvina, \emph{Questions in Geometric Group Theory} (2004), Question 8.2, PDF page 16, attributed to Shalen: \url{https://www.math.utah.edu/~bestvina/eprints/questions-updated.pdf}; manual transcription, accessed 2026-07-29.
\item Dunwoody, arXiv:1203.6019v3 (2022), \url{https://arxiv.org/abs/1203.6019v3}, the affirmative preprint claim with version-history and verification caveats recorded in the supplied summary.
\item Minasyan, \emph{Journal of Topology} (2016), cited for finitely generated counterexamples with property (FA) and finite arc stabilizers. No title, volume, pages, or theorem number is supplied.
\item Local registry \texttt{data/unsolvedmath/problems.json}, permanent ID $1100802$, source code catalogue entry, \emph{Questions in Geometric Group Theory --- Q 8.2}. The record labels this partially solved and PARTIAL-PROGRESS; the source-list background marks current-status and rights checks NEEDS\_REVIEW.
\end{itemize}
\end{document}
